\documentclass[11pt]{article}
\usepackage[margin=1in]{geometry}
\usepackage[T1]{fontenc}
\usepackage{xurl}
\setlength{\emergencystretch}{2em}
\title{Tactix: DSR checker-route gate and WV100 launch}
\author{Codex for Andrew}
\date{27 September 2026, 21:56 UTC}
\begin{document}
\maketitle

\paragraph{Objective and outcome.}
The immediate objective was to make the newly approved native-DSR verification
contract operational, then submit the full 100-instance winner-vs-Tactix
(WV100) Phase 2 gate. This increment is implemented and pushed as source commit
\texttt{0a906338}; pinned array \texttt{1731686\_0-99} is running. At the
report timestamp, 91/100 final records are independently checked and valid:
30 original-CNF SAT models and 61 complete UNSAT proof routes. Nine rows,
including \texttt{ncc\_none\_21015}, remain. The full Phase 2 gate has not
been accepted, and no Tier D job has been submitted.

\paragraph{Checker contract.}
Every native DSR proof now goes first through pinned no-option
\texttt{dsr-trim}. A witness-free LSR elaboration finishes with CakeLPR;
one containing an SR witness finishes with Trestle. The solver writes one
\texttt{c proof-route} declaration on an UNSAT proof-producing run based on
whether its proof writer actually emitted an SR witness. The production
runtime follows that declaration and scans the certificate only to reject a
mismatch. The independent checkers still see a clause-sequence-preserving
spelling of the original CNF. Startup PBP remains VeriPB then CakePB; explicit
DRAT/LRAT remain compatibility paths. No legal public X7 action or proof rule
was narrowed, and no row-specific checker setting was introduced. This change
replaces the earlier witness-free DRAT-trim route because the user explicitly
selected the general DSR-trim-first contract after the archived \texttt{ncc}
diagnostic; the current-source \texttt{ncc} certificate is a different proof,
so the archived timing cannot certify this campaign.

\paragraph{Correctness evidence.}
The checker/runtime/fuzz protocol suite passed 113 Python tests; Release and
Debug proof serializer tests and the current-source Debug \texttt{micro\_proofs}
suite passed against independent checkers. The previous repaired-source full
Debug sweep completed 188/188 but predates this route edit. The accepted
pre-route Phase 1 arrays checked Rank 1--50 at 50/50 (25 SAT models, 25 UNSAT
proofs, 7,553.134 solver seconds versus a strict 8,115.735-second target) and
B28+C5 at 33/33 (17 models, 16 proofs, 905.892 versus 952.263 seconds).
Latest-source proofless Rank finished 50/50 with all 25 SAT models checked;
proof-producing performance and proof checks under the new route still need
current-source confirmation. At 21:56 UTC, the 91 AM507 final records had
matching verdicts, complete independent checks and cleanup; 30 were SAT, 61
UNSAT. This is interim evidence only: the pinned scorer requires exactly 100
final records, 33 SAT models and 67 UNSAT proofs. Early UNSAT routes include
both CakeLPR and Trestle after \texttt{dsr-trim}.

\paragraph{Resource and provenance boundary.}
The submitted bundle freezes source \texttt{0a906338}, Release binary SHA-256
\texttt{cab44387ea47f3ed5d74cbf1a9bb702f321cf9164a3df791ef4b07d760c08656},
all 100 input hashes, and seven checker pins. The full manifest digest is
recorded in the AM507 launch report. The unthrottled 100-row array requests
\texttt{cpu-dense-hwulab-q},
12 CPUs and 32 GiB per task, a 1,200-second solve window and a 57,600-second
checker window without overrides. The older AM500 old-route bundle is frozen
and unsubmitted. These limits belong to WV100, not the eventual Tier D run.

\paragraph{Unresolved obligations.}
The exact two-stage exit definition is proposed in the AM508 acceptance
report. Stage A is the finite, independently checked preflight that authorizes
the full Tier D submission; Stage B scores that run and decides N3. This avoids
requiring Tier D evidence before Tier D can be launched. The proposal is not a
waiver or an accepted change to the plan.

First, finish AM507 and investigate any timeout or rejected proof generally.
Then reconfirm current-source Phase 1 and full Debug coverage, run fresh
Opus-owned default/random-policy fuzzing, establish public reach of the
direct-Proof-API D2F retained-deletion failure, and close broad N1/N2
proof/model and rule-action coverage. The CaDiCaL/Kissat clausal emitter map
is structurally complete, while the general live pseudo-Boolean arithmetic and
redundancy interface used by LymphoSAT/Satsuma remains a documented gap.
N3 still needs matched Rank attribution, checked WV2--89 at its 600-second
limit, Tier B/C parity and Tier D Kissat/VBS performance. N4 still needs its
scheduling parity ledger closed or explicitly accepted. An online-learning
readiness review must address adjustable inprocessing dose/order and
switchable branching and clause-deletion strategies with correctness and
overhead evidence. Only after the agreed gates are verified or explicitly
revised should the full 1,166-instance Tier D array be submitted at 5,000
seconds, four CPUs and 32 GiB per instance. Any relaxed gate must be reported
as unverified, not passed.

\paragraph{Primary evidence.}
\url{benchmarks/reports/2026-09-27-am506-dsr-trim-first-gate.md};
\url{benchmarks/reports/2026-09-27-am507-wv100-dsr-trim-first-launched.md};
\url{benchmarks/reports/2026-09-27-am508-goal-exit-contract-proposal.md};
\url{benchmarks/reports/2026-09-27-am502-repaired-phase1-terminal.md};
\url{benchmarks/reports/2026-09-27-am505-proofless-rank-terminal.md};
\url{benchmarks/reports/2026-09-27-am493-proof-rule-coverage-current-pin.md};
\url{benchmarks/reports/2026-09-25-am391-learning-readiness-current-source.md}.
\end{document}
